\documentclass[12pt, titlepage]{article}
%\usepackage[norsk]{babel}
\usepackage[utf8]{inputenc}
\usepackage{minted, enumerate}
\usepackage[colorlinks = true,
linkcolor = blue,
urlcolor = blue]{hyperref}


\begin{document}

\title{Emneinfo \\ Endringer til EPN 24/25}
\author{JJH}
\date{}
\maketitle


\section*{Faglig innhold}

\subsection*{Gammel}
Emnet gjev studenten ei grunnleggjande forståing for bruk av moderne digitale verkty og metodar for å løysa problem knytt til økonomi og leiing.  Me går mellom anna inn i fylgjande tema:

    Grunnleggjande bruk av programmering i økonomifag
    Omgrep som vert brukte i informasjonsteknogi
    Korleis ein analyserer store og fragmenterte datasett som stammar frå fleire kjelder
    Visualisering av data som verkemiddel i kommunikasjon og formidling (kundedata, rekneskapsdata, logistikk, osv)
    Forståing for datasikkerheit
    Kjelder til informasjon og data for analyse
    Korleis kunstig intelligens og maskinlæring kan nyttast i økonomifaga

\subsection*{NY}

hei

\section*{Læringsutbytte}
\subsection*{Gammel}

\subsubsection*{Kunnskap}

    Studenten skal kjenna høve og avgrensingar ved programmering og andre digitale verkty for handsaming og analyse av økonoiske data.
    Studenten skal kjenna til dei grensene som lovverket set for datahandsaming.

\subsubsection*{Ferdigheiter}

    Student skal kunna bruka programmering for å simulere, modellera og løyse problemstillingar i økonomiske fag.

\subsubsection*{Generell kompetanse}

    Studenten skal kunna kommunisera om dataverkty, -behov og -løysingar på tvers av faggrenser.
    Studenten skal kunna bruka simulering og dataanalyse som grunnlag for økonomiske vurderingar.
    Studenten skal kunna identifisera og vurdera informasjonsverdiar og sikkerheitsrisiko knytt til datahandsaming.
   
\subsection*{Ny}
hei


\section*{Læringsformer aktiviteter}
\subsection*{Gammel}
    Underviste økter med interaktive og studentaktive læringsformer (t.d. live koding, gruppearbeid)
    Praktiske øvingar og prosjekt under rettleiing med læringsassistentar og faglærar.

\subsection*{Ny}

Beholder i sin helhet.
Forelelesningstunge deler av emnet vurderes gjøres om til omvendt undervisning med video+quiz i forkant - praktisk arbeid i time



\end{document}
